\documentclass[11pt,a4paper]{article}

\usepackage[utf8]{inputenc}
\usepackage[T1]{fontenc}
\usepackage[spanish]{babel}
\usepackage[margin=2cm, top=2.6cm, bottom=1.8cm, footskip=0.6cm]{geometry}
\usepackage{graphicx}
\graphicspath{{../../01_Propuesta/Logos/}{./}}
\usepackage{fancyhdr}
\usepackage{xcolor}
\usepackage[hidelinks]{hyperref}
\usepackage{enumitem}
\usepackage{array}
\setlist{nosep}
\setlength{\parindent}{0pt}
\setlength{\parskip}{4pt}

\definecolor{verdeosc}{RGB}{26,107,60}
\definecolor{verdemid}{RGB}{46,139,87}
% Datos que faltan: se ven en rojo hasta completarlos.
\newcommand{\pend}[1]{\textcolor{red}{[PENDIENTE: #1]}}

\setlength{\headheight}{40pt}
\pagestyle{fancy}
\fancyhf{}
\fancyhead[L]{%
  \begin{tabular*}{\textwidth}{@{\extracolsep{\fill}}cccc@{}}
    \includegraphics[height=1.0cm]{logo_unisucre.png} &
    \includegraphics[height=1.0cm]{logo_sicvec2026.png} &
    \includegraphics[height=1.0cm]{Logo_insilico.jpeg} &
    \includegraphics[height=1.0cm]{Logo_Bio.png} \\
  \end{tabular*}%
}
\fancyfoot[C]{\footnotesize\itshape SICVEC 2026 contribuye a los Objetivos de Desarrollo Sostenible 12, 13, 15 y 17.}
\renewcommand{\headrulewidth}{1pt}
\renewcommand{\headrule}{\hbox to\headwidth{\color{verdeosc}\leaders\hrule height\headrulewidth\hfill}}

\begin{document}

\vspace{0.1cm}

\noindent
\textbf{Sincelejo, 2 de octubre de 2026}

\vspace{0.1cm}

\noindent
Señor\\
Esp. Carlos Meza Bertel\\
Director Ambiental y de Servicios Generales\\
Parque Comercial Guacarí\\
Sincelejo, Sucre, Colombia

\vspace{0.1cm}

\noindent
\textbf{\large Asunto: Solicitud de auditorio y espacios para SICVEC 2026}

\vspace{0.1cm}

Estimado señor Meza:

\vspace{0.1cm}

En nombre del Grupo de Investigación IN SILICO y del Departamento de Biología y Química de la Universidad de Sucre (UNISUCRE), me permito solicitar formalmente al Parque Comercial Guacarí el uso de su auditorio y de dos espacios complementarios para la realización del \textbf{Simposio Internacional de Ciencia Verde y Economía Circular (SICVEC 2026)}, los días \textbf{lunes 19 y martes 20 de octubre de 2026}.

El simposio reunirá a investigadores, estudiantes, profesionales y empresarios en torno a la economía circular, la química verde y las tecnologías sostenibles. Se realizará en modalidad híbrida, con conferencistas magistrales de México, Perú, Italia, Brasil y Colombia, y esperamos entre 120 y 150 asistentes presenciales. El Guacarí figura como sede del evento en el sitio web, el afiche y los materiales de difusión.

\vspace{0.1cm}

\noindent
\textbf{ESPACIOS SOLICITADOS}

\renewcommand{\arraystretch}{1.15}
\begin{tabular}{@{}>{\raggedright\arraybackslash}p{3.6cm}>{\raggedright\arraybackslash}p{4.4cm}>{\raggedright\arraybackslash}p{8.2cm}@{}}
\textbf{Espacio} & \textbf{Fecha y horario} & \textbf{Uso y requerimientos} \\
\hline
\textbf{Auditorio (sala de conferencias)} & Lunes 19 de octubre, 07:45--17:30\newline Martes 20 de octubre, 08:00--17:00 & Acto de apertura, conferencias magistrales, ponencias y clausura. Sonido, proyección y conexión a internet para la transmisión en línea. \\
\textbf{Espacio para la muestra empresarial} & Lunes 19 y martes 20 de octubre, durante la jornada & Stands de empresas y emprendimientos verdes, con mesa, sillas y punto eléctrico. \\
\textbf{Espacio para la sesión de pósteres} & Lunes 19 de octubre, 16:10--17:30 & Exhibición de pósteres científicos de 90~$\times$~120~cm; superficie o paneles para fijarlos y espacio de circulación para los asistentes. \\
\textbf{Acceso previo} & Domingo 18 de octubre, 4:00 p.\,m. & Montaje y prueba de los equipos audiovisuales y de la transmisión en línea. \\
\end{tabular}
\renewcommand{\arraystretch}{1}

\vspace{0.1cm}

El número definitivo de pósteres se conocerá el 15 de octubre, cuando concluya la evaluación de los resúmenes, y se lo comunicaremos junto con el número de stands de la muestra empresarial. La organización se compromete a entregar los espacios en las mismas condiciones en que los reciba y a acatar las normas de uso del Parque Comercial.

El Coordinador Logístico del simposio, Sebastián Vargas, queda a su disposición para concertar una visita a los espacios y definir con su equipo los detalles de montaje, aforo y horarios.

\vspace{0.1cm}

Agradezco de antemano su apoyo y quedo atento a su respuesta. Contar con el Parque Comercial Guacarí como sede es un respaldo muy valioso para SICVEC 2026 y para la Universidad de Sucre.

\vspace{0.15cm}

Atentamente,

\vspace{0.15cm}

\textbf{Dr. Aldo F. Combariza}\\
Coordinador General, SICVEC 2026\\
Grupo de Investigación IN SILICO · Departamento de Biología y Química\\
Facultad de Educación y Ciencias · Universidad de Sucre\\
\href{mailto:aldo.combariza@unisucre.edu.co}{aldo.combariza@unisucre.edu.co}

\vspace{0.1cm}

\end{document}
